% Part: applied-modal-logics
% Chapter: epistemic-logic
% Section: truth-at-w

\documentclass[../../../include/open-logic-section]{subfiles}

\begin{document}

\olfileid{aml}{el}{trw}

\olsection{在世界中为真}

与正规模态逻辑一样，每个认知模型都决定了其中哪些公式在哪些世界中为真。
我们使用相同的记号“模型 $\mModel{M}$ 使公式~$!A$ 在世界~$w$ 为真”来表达关系语义的基本概念。
这一关系是归纳定义的，并且对所有非模态算子而言，与正规模态逻辑的情形完全相同。

\begin{defn}\ollabel{defn:mmodels}
  公式~$!A$ 在模型~$\mModel M$ 的世界~$w$ 处的\emph{真}，记作：
  $\mSat{M}{!A}[w]$，归纳定义如下：
  \begin{enumerate}
  \tagitem{prvFalse}{\indcase{!A}{\lfalse}{永不 $\mSat{M}{\lfalse}[w]$}.}{}
  \tagitem{prvTrue}{\indcase{!A}{\ltrue}{恒有 $\mSat{M}{\ltrue}[w]$}.}{}
  \item $\mSat{M}{p}[w]$ 当且仅当 $w \in V(p)$
  \tagitem{prvNot}{\indcase{!A}{\lnot !B}{$\mSat{M}{\indfrm}[w]$ 当且仅当
    $\mSat/{M}{!B}[w]$}.}{}
  \tagitem{prvAnd}{\indcase{!A}{(!B \land !C)}{$\mSat{M}{\indfrm}[w]$ 当且仅当
    $\mSat{M}{!B}[w]$ 且 $\mSat{M}{!C}[w]$}.}{}
  \tagitem{prvOr}{\indcase{!A}{(!B \lor !C)}{$\mSat{M}{\indfrm}[w]$ 当且仅当
    $\mSat{M}{!B}[w]$ 或 $\mSat{M}{!C}[w]$（或两者皆成立）}.}{}
  \tagitem{prvIf}{\indcase{!A}{(!B \lif !C)}{$\mSat{M}{\indfrm}[w]$ 当且仅当
    $\mSat/{M}{!B}[w]$ 或 $\mSat{M}{!C}[w]$}.}{}
  \tagitem{prvIff}{\indcase{!A}{(!B \liff !C)}{$\mSat{M}{\indfrm}[w]$ 当且仅当
    要么 $\mSat{M}{!B}[w]$ 与 $\mSat{M}{!C}[w]$ 均成立，要么 $\mSat{M}{!B}[w]$ 与 $\mSat{M}{!C}[w]$ 均不成立}.}{}
  \item\ollabel{defn:sub:mmodels-box}
    \indcase{!A}{\Knows_a !B}{$\mSat{M}{\indfrm}[w]$ 当且仅当
    对所有满足 $R_a ww'$ 的 $w' \in W$，均有 $\mSat{M}{!B}[w']$}
  \end{enumerate}
\end{defn}

不过，在此处我们需要考虑对可达关系的限制。
毕竟，根据条款~\olref{defn:sub:mmodels-box}，当不存在满足 $R_a ww'$ 的 $w'$ 时，公式~$\Knows_a !B$ 在世界~$w$ 就为真。
这与正规模态逻辑中的条款相同；当一个世界没有后继时，所有 $\Box$-公式在那里都空真。
如果我们将 $\Knows$ 视为代表\emph{知识}，这似乎极度违反直觉。
毕竟，我们通常认为，\emph{没有}任何情况能使一个智能体同时知道 $!A$ 和 $\lnot !A$。

一种解决方案是确保认知逻辑中的可达关系总是\emph{自反的}。这大致对应于这样的想法：现实世界与智能体的信息是一致的。事实上，认知逻辑通常使用 S5 系统，但其他人可能会使用更弱的系统，这取决于他们究竟希望 $\Knows_a$ 关系代表什么。

\begin{figure}
  \begin{center}
    \begin{tikzpicture}[modal]
      \node[world] (w1) [label={[align=right]right:\mTrue{p}\\ \mFalse{q}}]{$w_1$}; 
      \node[world] (w2) [label={[align=right]right:\mTrue{p}\\ \mTrue{q}},
        above right=of w1]{$w_2$}; 
      \node[world] (w3) [label={[align=right]right:\mFalse{p}\\ \mFalse{q}},
        below right=of w1] {$w_3$};
      \draw[<->] (w1) to node {$a$} (w2);
      \draw[->,loop left] (w1) to node {$a,b$} (w1);
      \draw[<->] (w1) to [swap] node {$b$} (w3);
      \draw[->,loop above] (w2) to node {$a,b$} (w2) ;
      \draw[->,loop below] (w3) to node {$a,b$} (w3) ;
    \end{tikzpicture}
  \end{center}
  \caption{一个简单的认知模型}
  \ollabel{fig:simple}
\end{figure}

\begin{prob}
  考虑下列哪些公式在 \olref[aml][el][trw]{fig:simple} 中成立：
  \begin{enumerate}
    \item $\mSat{M}{\lnot q}[w_1]$;
    \item $\mSat{M}{\Knows_a \lnot q}[w_1]$;
    \item $\mSat{M}{\Knows_b \lnot q}[w_1]$;
    \item $\mSat{M}{\Knows_b q \lor \Knows_b \lnot q}[w_2]$;
    \item $\mSat{M}{\Knows_a( \Knows_b q \lor \Knows_b \lnot q)}[w_2]$;
    \item $\mSat{M}{\EKnows_{\{a,b\}} \lnot q}[w_3]$;
    \end{enumerate}
\end{prob}

既然我们已经给出了公式在世界中为真的基本定义，正规模态逻辑中的其他语义概念，例如模态有效性和语义后承，就可以直接推广过来，应用于这种关于模态算子解释的新思考方式。

现在我们也可以为公共知识算子~$\CKnows_G$ 给出真值条件。回顾 \olref[sfr][rel][ops]{sec}，关系~$R$ 的\emph{传递闭包}~$R^+$ 定义为
\begin{align*}
  R^+ &= \bigcup_{ n \in \mathbb{N}} R^n, 
\intertext{其中}
  R^0 & = R \text{ 且}\\ 
  R^{n+1} & = \Setabs{\tuple{x, z}}{\exists y (R^n xy \land Ryz)}.
\end{align*}
那么，若 $G$ 是一个智能体群，我们定义 $R_G = ( \bigcup_{b
\in G} R_b )^+$ 为所有智能体可达关系的并集的传递闭包。

\begin{defn}
若 $G' \subseteq G$，我们令 $\mSat{M}{\CKnows_{G'} !A}[ w]$ 当且仅当对于每一个满足 $R_{G'} w w'$ 的 $w'$，均有 $\mSat{M}{!A}[w']$。
\end{defn}

\end{document}